\section{Experimental Setup}
\label{sec:experimental_setup}

% Follows G_paper_and_latex/paper_outline.md (Section 4) and
% E_research_design/experiment_matrix.md.
% All datasets, parameters, and baseline versions must be empirically verified.

\subsection{Datasets and Annotation Protocols}
\label{sec:exp_datasets}
% TODO: Official dataset name, release version, image counts, class distribution,
% split methodology (train/val/test) avoiding scene/camera leakage.
% Reference: templates/dataset_card.md and E_research_design/threats_to_validity.md (Section 1)
\noindent [TODO: Document official benchmark datasets (e.g., NVIDIA PhysicalAI-SmartSpaces and auxiliary candidate PhysicalAI-Spatial-Intelligence-Warehouse). Detail image resolutions, verified sample counts, object class definitions, annotation schemas, and train/val/test split protocols ensuring zero camera/scene leakage. See \texttt{E\_research\_design/threats\_to\_validity.md} and \texttt{templates/dataset\_card.md}].

\subsection{Data Preprocessing and Augmentation}
\label{sec:exp_preprocessing}
% TODO: Input resizing (e.g., 640x640), normalization, augmentations (Mosaic, MixUp, HSV jitter),
% and label format conversion.
\noindent [TODO: Detail data preprocessing procedures, image normalization, resolution scaling, and data augmentation policies applied during training and evaluation].

\subsection{Baseline Architectures and Experimental Cases}
\label{sec:exp_baselines}
% Follows yolo_mamba_research_cases.md and E_research_design/experiment_matrix.md
To establish rigorous, unconfounded comparisons, our evaluation is organized into phased experimental cases centered on \textbf{YOLO11s}:
\begin{itemize}
    \item \textbf{Case A (Reference Baseline):} Unmodified standard YOLO11s trained under identical hyperparameter and data configurations.
    \item \textbf{Case B (Primary Novelty / First Milestone):} Direct replacement of the deepest \texttt{C2PSA} attention block at $P_5$ (stride 32: $20 \times 20$) with a shape-compatible \texttt{SS2D} block.
    \item \textbf{Case C (Dual-Scale SSM):} Integration of \texttt{SS2D} blocks across both $P_4$ (stride 16: $40 \times 40$) and $P_5$ to test whether mid-level state-space modeling justifies additional latency.
    \item \textbf{Case D (Supplemental SSM):} Co-existence of \texttt{C2PSA} and \texttt{SS2D} at $P_5$ to isolate whether attention and SSM features are complementary or redundant.
    \item \textbf{Case E (Context-Gated Neck):} Evaluation of $P_5 \to P_3$ contextual gating against the original PAN/FPN neck on the best-performing backbone variant.
\end{itemize}

\subsection{Implementation and Training Details}
\label{sec:exp_implementation}
% TODO: Framework versions (PyTorch, CUDA), optimizer (AdamW/SGD), initial lr,
% cosine scheduler, weight decay, batch size, epoch budget, early stopping, random seeds.
% Reference: F_literature_and_evidence/reproducibility_manifest.md
\noindent [TODO: Specify complete software and hardware training environments: PyTorch version, CUDA runtime, GPU hardware, optimizer parameters, learning rate schedules, batch sizes, epoch limits, and fixed random seeds. See \texttt{F\_literature\_and\_evidence/reproducibility\_manifest.md}].

\subsection{Evaluation Metrics}
\label{sec:exp_metrics}
% TODO: Detection accuracy: mAP@0.5, mAP@0.5:0.95, Precision, Recall, F1-score.
% Efficiency metrics: Parameter count (M), FLOPs/MACs (G), Model size (MB), latency (ms), FPS, peak RAM (MB).
\noindent [TODO: Provide mathematical definitions for all evaluation metrics: Precision, Recall, mean Average Precision ($\text{mAP}@0.5$ and $\text{mAP}@0.5:0.95$), parameter counts, computational complexity ($\text{GFLOPs}$), peak memory, and frame-rate ($\text{FPS}$)].

\subsection{Edge Deployment and Benchmarking Protocol}
\label{sec:exp_edge_protocol}
% TODO: Edge hardware specifications (Jetson Orin Nano, Raspberry Pi 4B/5), runtime engines
% (TensorRT, ONNX Runtime), precision modes (FP32, FP16, INT8), batch=1, warm-up iterations,
% timing boundaries (end-to-end vs model-only).
% Reference: E_research_design/deployment_plan.md
\noindent [TODO: Describe the standardized edge hardware benchmarking protocol: target edge device(s), runtime backends (e.g., ONNX Runtime, TensorRT), precision formats, batch size ($=1$), warm-up iterations, repetition counts, synchronization rules, and explicit timing boundaries (model-only vs.\ end-to-end pipeline). Reference \texttt{E\_research\_design/deployment\_plan.md}].
